\documentclass[10pt,a4paper]{amsart}
\usepackage[margin=19mm]{geometry}
\usepackage{amsmath,amssymb,mathtools}
\usepackage[scr=rsfs,cal=euler]{mathalfa}
\usepackage{xfrac}
\usepackage[unicode,hidelinks,bookmarks=false]{hyperref}
\newcommand{\ie}{i.e.\ }
\newcommand{\cf}{cf.\ }
\newcommand{\resp}{respectively\ }
\newcommand{\Subsection}{\subsection}
\providecommand{\nobreakdash}{\nobreak\hbox{-}}
\setlength{\emergencystretch}{2em}
\multlinegap0pt
\usepackage{amssymb}
\usepackage{xcolor}
\newtheorem{lemma}{Lemma}
\newtheorem{proposition}{Proposition}
\newcommand{\C}{\mathcal C}
\newcommand{\DD}{\mathbf D}
\newcommand{\E}{\mathcal E}
\newcommand{\F}{\mathcal F}
\renewcommand{\H}{\mathcal H}
\renewcommand{\O}{\mathcal O}
\newcommand{\RHom}{{\bf R} \H om}
\DeclareMathOperator{\depth}{depth}
\newcommand{\sExt}{\E xt} %sheaf Ext
\DeclareMathOperator{\supp}{supp}
\newcommand{\xto}{\xrightarrow} %use \xto^{\sim} to add texts/symbols on arrows
\title{Injectivity and vanishing for the Du Bois complexes of isolated singularities}
\author{Mihnea Popa, Wanchun Shen,, Anh Duc Vo}
\date{Source: arXiv:2409.18019v2 (CC BY 4.0)}
\begin{document}
\maketitle

\noindent\emph{Source and licence.} Injectivity and vanishing for the Du Bois complexes of isolated singularities. Version arXiv:2409.18019v2 (CC BY 4.0), distributed by arXiv under the Creative Commons Attribution 4.0 International licence, http://creativecommons.org/licenses/by/4.0/, as recorded in the arXiv licence metadata for this exact version and shown on the abstract page https://arxiv.org/abs/2409.18019v2. Full licence notice: \texttt{licenses/arxiv-2409.18019-CC-BY-4.0.txt}.\par\smallskip

\noindent\textit{Editorial scope.} Counted unit: the article's proposition prop:derived (source0.tex lines 950--959). The article's complete original proof of it is delivered with it (source0.tex lines 960--986, one \texttt{proof} environment of 27 lines). No mathematics is written, repaired or re-derived: the statement and the proof are the article's own lines, in the article's own notation, and no retained line is substituted or re-flowed.  Retained uncounted as the source-local context the proof needs: lines 671--673; lines 702--705.  Nothing was omitted from any retained span, so the package carries no omission record.  No retained line contains a citation, so the wrapper carries no bibliography and no bibliography entry.  No retained line contains a figure environment, an image-inclusion command or drawing code, so no image is delivered and the package declares no asset dependency.  Editorial formatting only, in addition: an AMS \texttt{amsart} preamble that restates the article's own declarations and macros used by the retained lines, namely the environments \texttt{lemma}, \texttt{proposition} on separate counters, together with the standard packages those lines need.  The retained environments print in this document as Lemma 1, Proposition 1 in order of appearance, whereas the article prints the same items with the numbering of its own full document.  The complete native source of the article as distributed in the arXiv e-print for this exact version is bundled unchanged as \texttt{sources/165869/source0.tex}.
\begin{equation}\label{eqn:depth}
\depth(\C) = \min~  \{i ~|~\sExt^{-i}_{\O_X} (\C,\omega_X^{\bullet})\neq 0\}.
\end{equation}

\begin{lemma}\label{lem:vanishing-of-Ext}
Let $\F$ be a coherent sheaf on $X$. Then 
\[\sExt^i_{\O_X}(\F, \omega_X^{\bullet})=0 \text{ \,\,{\rm for} \,\,} i< -\dim {\rm Supp}( \F).\]
\end{lemma}

\begin{proposition}\label{prop:derived}
Let $A^\bullet$ be an object in $\DD^b_{\rm coh} (X)$ such that 
\begin{enumerate}
\item $A^\bullet$ has cohomology in non-negative degrees.
%, and $\H^0 A^\bullet \neq 0$.
\item The support of all $\H^i A^\bullet$ with $i > 0$ is contained in a closed subset of $X$ of dimension $s$.
\end{enumerate}
Then we have
\[\H^i A^\bullet  = 0 \quad \text{ {\rm for} } \quad 0<i< {\rm min} \{\depth \H^0 A^\bullet, \depth A^\bullet + 1\} - s -1.\]
\end{proposition}

\begin{proof}
The first assumption implies that there is a natural morphism $\H^0 A^\bullet  \to A^\bullet$, and taking its Grothendieck dual 
leads to
\begin{equation} \label{eqn:cone of non-DB locus}
\RHom_{\O_X}(A^\bullet,\omega_X^{\bullet}) \xto{\varphi} \RHom_{\O_X}(\H^0 A^\bullet,\omega_X^{\bullet}) \to C^{\bullet}\xto{+1}, 
\end{equation}
where $C^{\bullet}$ is the cone of the morphism $\varphi$. For each $q$, we obtain exact sequences
\[\sExt^q_{\O_X}(\H^0 A^\bullet,\omega_X^{\bullet})\to \H^q C^{\bullet}\to  \sExt^{q+1}_{\O_X}(A^\bullet,\omega_X^{\bullet}). \]
Using ($\ref{eqn:depth}$), the first term vanishes for $q>-\depth \H^0 A^\bullet$, and the third for  $q> -\depth A^\bullet - 1$. We conclude that 
\[\H^q C^{\bullet}=0\text{ for } q> - m,\]
where $m : = {\rm min} \{\depth \H^0 A^\bullet, \depth A^\bullet\}$.

Next, applying $\RHom_{\O_X}(-, \omega_X^{\bullet})$ to (\ref{eqn:cone of non-DB locus}), we obtain an exact triangle 
\[\RHom_{\O_X}(C^{\bullet},\omega_X^\bullet)\to \H^0 A^\bullet  \to A^\bullet \xto{+1}.\]
It follows that for $i>0$,
\[\H^i A^\bullet \simeq \sExt^{i+1}_{\O_X}(C^\bullet, \omega_X^{\bullet}).\]
Consider now the spectral sequence 
\[E^{p,q}_2 = \sExt^p_{\O_X}(\H^q C^{\bullet},\omega_X^\bullet) \implies \sExt^{p-q}_{\O_X}(C^{\bullet},\omega_X^\bullet).\]
We have:
\begin{itemize}
    \item $\H^qC^{\bullet}=0$ for $q> - m$, as noted above;
    \item $\sExt^p_{\O_X}(\H^q C^{\bullet},\omega_X^{\bullet})=0$ for $p<-\dim \supp \H^q C^{\bullet}$, by Lemma \ref{lem:vanishing-of-Ext}. In particular, this holds for $p<-s$, since by definition $C^{\bullet}$ is supported on the locus where $\H^0 A^\bullet$ and $A^\bullet$ are not 
quasi-isomorphic, which has dimension $s$.
\end{itemize}
Combining these facts, we see that $\sExt^i_{\O_X}(C^{\bullet},\omega_X^{\bullet})=0$ for $i< m - s$. Thus
\[\H^i A^\bullet = 0 \,\,\,\,\,\,{\rm for} \,\,\,\,0<i< m - s -1.\]
\end{proof}

\end{document}
